% A short, real beautybook example. Build with the supplied menu script.
\documentclass[lang=cn,zihao=-4,a4paper,fontset=none]{beautybook}
\beautybookstyle={cover-choose=cn,math-font=plain,sidebar=off}
\usepackage{stys/beautybook-cnsettings}
\usepackage{stys/beautybook-liquidglass}
\chapimage{c10}
\partimage{c11}
\setlength\parindent{2em}
\begin{document}
\mainmatter
\chapter{玻璃材质已经接好}
本例沿用这份工程原来的文类、字体和定理环境。编号、可选标题与引用的写法保持不变。
\section{原有盒子直接使用}
\begin{defi}[第一套定理]
旧写法 \texttt{defi} 自动获得玻璃表面，定理正文与数学公式仍由 \LaTeX{} 排版。
\[ (a+b)^2=a^2+2ab+b^2. \]
\end{defi}
\begin{problem}[][第二套定理][demo:problem]
问题编号与旧的彩色标题装饰保留。这里的背景来源为随包默认素材。
\end{problem}
\begin{Definition}[note={第三套定理},label={demo:definition}]
大写的 \texttt{Definition} 来自原模板的 keytheorems 系统，同样使用玻璃表面。
\end{Definition}
\begin{fancybox}
原来的 \texttt{fancybox} 环境也可以照常调用。
\end{fancybox}
\begin{tcolorbox}[title={普通 tcolorbox},arc=3mm,breakable]
正文区的普通盒子在选项读取完成后加入材质，因此旧背景填充不会再盖住玻璃层。
\end{tcolorbox}
\begin{glassbox}[title={逐盒选择深色},lg theme=dark]
显式使用 \texttt{glassbox} 时，当前盒子的深浅色与参数设置优先。
\end{glassbox}
\section{长内容自动跨页}
下方使用原来的 \texttt{definition} 环境，演示同一逻辑盒子的连续背景与分页。
\begin{definition}[][跨页定义][demo:long]
\newcount\demoParagraph\demoParagraph=0
\loop\ifnum\demoParagraph<28
\advance\demoParagraph by1
\textbf{第\the\demoParagraph 段。} 背景只在一个逻辑盒子开始时确定，跨页后继续使用同一来源；正文、公式与编号始终清晰。\quad $a^2+b^2=c^2$。\par\medskip
\repeat
\end{definition}
\section{小白用户只需改这两行}
\begin{verbatim}
\BeautyGlassBackground{liquidglass-assets/default.png}
\BeautyGlassSetup{lg reference}
\end{verbatim}
把设置放在主文件导言区加载 \texttt{stys/beautybook-liquidglass} 之后即可。\par
问题 \ref{demo:problem} 与定义 \ref{demo:definition} 的标签仍按原模板引用。
\end{document}
